\documentclass[journal,twoside]{IEEEtran}
\usepackage{amsmath,amssymb,amsfonts,amsthm}
\usepackage{graphicx}
\usepackage{xcolor}
\usepackage{booktabs}
\usepackage{algorithm}
\usepackage{algpseudocode}
\usepackage{cite}
\usepackage{url}
\usepackage{hyperref}

\hypersetup{colorlinks=true,linkcolor=blue,citecolor=blue,urlcolor=blue}

\newtheorem{theorem}{Theorem}
\newtheorem{assumption}{Assumption}
\newtheorem{remark}{Remark}
\newtheorem{definition}{Definition}

\newcommand{\norm}[1]{\left\|#1\right\|}
\newcommand{\R}{\mathbb{R}}
\newcommand{\ehat}{\hat{e}}
\newcommand{\ahat}{\hat{a}}

\begin{document}

\title{Deep-Learned Kernel Robust Moving Horizon Estimation for
Cyber-Resilient Closed-Loop Artificial Pancreas Digital Twins
Under Composite FDI/DoS/Replay Attacks}

\author{Subrahmanyam~Tanala,~\IEEEmembership{Member,~IEEE}%
\thanks{S.~Tanala is with the Department of Electrical and Electronics Engineering,
ANITS, Visakhapatnam 530003, India, and Woosong University, Daejeon 34606,
South Korea (e-mail: subrahmanyam.tanala@anits.edu.in).}%
\thanks{Code: \url{https://github.com/subrahmanyamtanala-rgb/dlk-rmhe-ap-cdt}}}

\markboth{IEEE Transactions on Biomedical Engineering}%
{Tanala: DLK-RMHE for Cyber-Resilient Artificial Pancreas Digital Twins}

\maketitle

\begin{abstract}
Clinical digital twins (CDTs) for artificial pancreas (AP) systems represent a
transformative paradigm in closed-loop Type~1 diabetes mellitus (T1DM) management.
The wireless connectivity of continuous glucose monitor (CGM) sensors, insulin
infusion pumps, and cloud-based model predictive controllers introduces a critical
cyber-physical attack surface. This paper proposes the \emph{Deep-Learned Kernel
Robust Moving Horizon Estimation} (DLK-RMHE) framework, integrating a
contrastively-trained deep kernel network within a sparse $\ell_1$-regularized MHE
to achieve provably resilient state estimation under simultaneous False Data
Injection (FDI), Denial-of-Service (DoS), and replay attacks on the CGM channel.
A dual-mode architecture handles DoS via Hovorka model prediction fallback and
FDI/replay via kernel-weighted measurement rejection. A complete Lyapunov ISS proof
uses the BPDN stable recovery framework with explicit derivation of constants
$C_1{=}2.0$, recovery bound $\|\ehat_k{-}a_k\|\leq 3.32$\,mmol/L, and ISS gains
$\gamma_a{=}20.9$, $\gamma_\varepsilon{=}2.20$. Assumption~A2 is justified via
Rademacher complexity giving $\varepsilon_{\max}\leq 0.797$. The deep kernel (four
FC layers: 128-64-32-16, ReLU, AdamW $\mathrm{lr}{=}10^{-3}$, 200~epochs,
50,000 pairs, PyTorch~2.1/RTX~3090) is fully specified for reproducibility.
Validation on the FDA-accepted UVa/Padova T1DM simulator (10 subjects,
$n{=}10$ Monte Carlo trials, 300 runs) achieves MAGE $4.2{\pm}1.1$\,mg/dL
(95\,\% CI: [3.9, 4.5]) under composite attack---84.2\,\% below standard
EKF---with TIR $76.4{\pm}5.1$\,\% and AUROC $0.962{\pm}0.015$. A four-variant
ablation study isolates each component contribution.
\end{abstract}

\section*{Notation}
\begin{IEEEdescription}[\IEEEusemathlabelsep\IEEEsetlabelwidth{$\varphi_\theta$}]
  \item[$x_k\in\R^7$] Hovorka state vector at step $k$
  \item[$\hat x_k$] DLK-RMHE state estimate
  \item[$e_k=x_k-\hat x_k$] Estimation error
  \item[$\ehat_k$] MHE slack variable (attack magnitude estimate)
  \item[$a_k$] True attack signal (FDI or replay)
  \item[$y_k,\,\tilde y_k$] Clean and corrupted CGM measurement
  \item[$\varepsilon_k$] Kernel approximation error $\|\varphi_\theta(z_k)-\varphi^*(z_k)\|$
  \item[$s_k$] Kernel anomaly score
  \item[$\rho_k$] Kernel-based measurement weight $\in(0,1]$
  \item[$\varphi_\theta$] Deep kernel embedding network, $\R^{24}\to\R^{16}$
  \item[$c_\mathrm{nom}$] Nominal cluster centroid in embedding space
  \item[$\gamma_\mathrm{th}$] Detection threshold ($=3.84$)
  \item[$\delta_d$] Uniform detectability index ($=0.18$)
  \item[$\delta_s$] RIP constant of CGM measurement operator ($=0$)
  \item[$\gamma_a,\gamma_\varepsilon$] ISS $\mathcal{L}_\infty$ gain functions
  \item[$P_k$] MHE arrival cost (covariance) matrix
  \item[$Q_w,R_v$] Process and measurement noise covariances
  \item[$N_w$] MHE window length ($=12$ steps, 60\,min)
  \item[$\lambda$] $\ell_1$ penalty weight ($=0.8$)
\end{IEEEdescription}


\begin{IEEEkeywords}
Artificial pancreas, clinical digital twin, cyber-physical security,
deep kernel learning, false data injection, moving horizon estimation,
Type~1 diabetes mellitus.
\end{IEEEkeywords}

% ─────────────────────────────────────────────────────────────────────────────
\section{Introduction}\label{sec:intro}
% ─────────────────────────────────────────────────────────────────────────────

\IEEEPARstart{T}{ype}~1 diabetes mellitus affects 8.4 million individuals worldwide,
requiring continuous exogenous insulin to compensate for total pancreatic beta-cell
loss~\cite{idf2021}. The artificial pancreas (AP)---integrating a CGM sensor,
insulin pump, and closed-loop algorithm---achieves sustained Time-in-Range (TIR)
improvements with systems such as the Medtronic MiniMed~780G and iLet Bionic
Pancreas~\cite{boughton2022,bhargava2022}. Convergence of AP technology with
clinical digital twins (CDTs) opens new dimensions in personalized T1DM
management~\cite{grieves2014,laubenbacher2021}.

The Bluetooth Low Energy (BLE) and WiFi connectivity of CGM-pump-cloud architectures
introduces a severe cybersecurity vulnerability. Adversaries can execute:
(i)~\emph{FDI}---injecting fabricated glucose readings to trigger dangerous insulin
doses; (ii)~\emph{DoS}---blocking CGM transmissions; and (iii)~\emph{replay}---
substituting live readings with historical normoglycemic records to mask acute
hypoglycemia~\cite{sevil2020,ransford2012}. Commercial AP state estimation via EKF
or nominal MHE offers no adversarial robustness guarantees~\cite{patek2012,chakrabarty2018}.

This paper addresses this gap with five contributions:
(1)~DLK-RMHE: complete architecture and training specification
(Tables~\ref{tab:arch}--\ref{tab:hyper});
(2)~dual-mode graceful degradation with formal switching logic
(Algorithm~\ref{alg:dlkrmhe});
(3)~complete ISS proof via BPDN stable recovery
(equations~(\ref{eq:B1})--(\ref{eq:B5}));
(4)~Rademacher complexity derivation for Assumption~A2
(equations~(\ref{eq:A1})--(\ref{eq:A3}));
(5)~four-variant ablation study with 300 Monte Carlo runs
(Table~\ref{tab:ablation}).

% ─────────────────────────────────────────────────────────────────────────────
\section{Related Work}\label{sec:related}
% ─────────────────────────────────────────────────────────────────────────────

\subsection{Artificial Pancreas Control and State Estimation}
MPC with the Hovorka~\cite{hovorka2004} and Bergman~\cite{bergman1981} models is
the dominant AP paradigm~\cite{magni2008}. EKF remains the commercial standard,
though Gaussian assumptions fail under adversarial inputs~\cite{patek2012}. Adaptive
MHE~\cite{chakrabarty2018,messori2019}, zone MPC~\cite{dassau2013}, and
physics-informed neural glucose predictors~\cite{shi2023,rodriguez2023} extend the
state-of-the-art. Safety-critical AP control with formal guarantees is addressed
in~\cite{farina2021}.

\subsection{Cybersecurity of Medical CPS}
Landmark vulnerabilities were shown for pacemakers~\cite{halperin2008} and insulin
pumps~\cite{radcliffe2011}. Formal AP threat models~\cite{sevil2020,li2020}, SPC
anomaly detection~\cite{datta2019}, one-class SVM~\cite{sevil2021}, deep learning
intrusion detection~\cite{xiao2022}, federated anomaly detection~\cite{ali2023},
stealthy FDI~\cite{mosquera2022}, watermarking~\cite{askari2021}, BLE
vulnerabilities~\cite{hei2010}, and formal closed-loop AP security
analysis~\cite{kim2023apsec} span the attack-defense spectrum. None provide provably
stable estimation under composite simultaneous attacks.

\subsection{Digital Twin Security}
Security frameworks for digital twins address FDI~\cite{mo2015},
blockchain integrity~\cite{xu2022}, twin synchronization robustness~\cite{nguyen2021},
federated privacy~\cite{li2021fedbn}, homomorphic cardiac
CDTs~\cite{raissi2019}, and secure
aggregation~\cite{bonawitz2017}.

\subsection{Robust and Learning-Enhanced Estimation}
Robust MHE with $\ell_1$ rejection and ISS guarantees is established for general
nonlinear systems~\cite{rao2003,alessandri2008}. Deep kernel learning for CPS
estimation~\cite{lederer2021,chen2023} and GP safety-critical
estimators~\cite{umlauft2018} motivate the present kernel-within-MHE approach.
Kernel dynamical system embedding~\cite{scholkopf1997} and Koopman
operators~\cite{korda2018} provide theoretical foundations. The dual
Kalman~\cite{wan2001} underpins the joint state-parameter estimation in the
arrival cost formulation.

% ─────────────────────────────────────────────────────────────────────────────
\section{System Model and Threat Formulation}\label{sec:system}
% ─────────────────────────────────────────────────────────────────────────────

\subsection{Hovorka Pharmacokinetic Model}
The seven-dimensional state $x(t)=[G,Q_1,Q_2,I,x_1,x_2,x_3]^\top\in\R^7$
follows the Hovorka ODEs~\cite{hovorka2004}:
\begin{align}
  \dot Q_1 &= -[F_{01}+x_1]Q_1+k_{12}Q_2-F_R+\mathrm{EGP}_0[1-x_3]+U_G,\\
  \dot Q_2 &= x_1Q_1-[k_{12}+x_2]Q_2,\\
  \dot I   &= -k_eI+u/V_I,\\
  \dot x_i &= -k_{a,i}x_i+k_{b,i}I,\quad i=1,2,3.
\end{align}
Fourth-order Runge-Kutta discretization at $T_s{=}5$\,min gives
$x_{k+1}=f(x_k,u_k,d_k)+w_k$, $y_k=Cx_k+v_k$, $C=[1\,0\,\cdots\,0]$,
$\sigma_v=0.07G+2$\,mg/dL (Dexcom G6~\cite{dexcom2021}).

\subsection{AP-CDT Architecture}
The CDT-augmented AP (Fig.~\ref{fig:arch}) comprises Physical, Communication
(BLE/WiFi with attack injection points), Digital Twin, and Control (MPC) layers.
\begin{figure}[!t]
  \centering
  \includegraphics[width=\columnwidth]{fig1_architecture}
  \caption{AP-CDT four-layer architecture with DLK-RMHE estimator and CGM attack injection points.}
  \label{fig:arch}
\end{figure}

\subsection{Composite Threat Model}
\begin{definition}[FDI]
$\tilde y_k=y_k+a_k^{\mathrm{FDI}}$, $\|a_k^{\mathrm{FDI}}\|\leq\bar a{=}2.5$\,mmol/L.
\end{definition}
\begin{definition}[DoS]
Bernoulli dropout; CGM blocked for $\ell\leq\ell_{\max}{=}9$ steps (45\,min).
\end{definition}
\begin{definition}[Replay]
$\tilde y_k=y_{k-\tau}$, $\tau\in\{6,\ldots,24\}$ steps.
\end{definition}
\begin{definition}[Composite]
Simultaneous activation of any subset of \{FDI, DoS, Replay\}.
\end{definition}

% ─────────────────────────────────────────────────────────────────────────────
\section{DLK-RMHE Framework}\label{sec:framework}
% ─────────────────────────────────────────────────────────────────────────────

\subsection{Deep Kernel Network}

Input: $z_k=[\Delta y_{k-N_w:k}^\top,u_{k-N_w:k}^\top]^\top\in\R^{24}$
($N_w{=}12$). Architecture in Table~\ref{tab:arch}; kernel:
$k_\theta(z,z')=\exp(-\|\varphi_\theta(z)-\varphi_\theta(z')\|^2/(2\sigma^2))$,
$\sigma^2{=}1.0$. Training: 50,000 labeled pairs, contrastive loss
(margin $m{=}2$), hyperparameters in Table~\ref{tab:hyper}.
Attack score and threshold:
$s_k=\|\varphi_\theta(z_k)-c_{\mathrm{nom}}\|^2$; attack if $s_k>\gamma_{\mathrm{th}}{=}3.84$.

\begin{table}[!t]
  \caption{Deep Kernel Network Architecture}
  \label{tab:arch}
  \centering
  \begin{tabular}{lrrrrrr}
    \toprule
    Layer & In & Out & Act. & Drop & BN & Params\\
    \midrule
    FC-1 & 24  & 128 & ReLU & 0.20 & Y & 3,200\\
    FC-2 & 128 & 64  & ReLU & 0.20 & Y & 8,256\\
    FC-3 & 64  & 32  & ReLU & 0.20 & N & 2,080\\
    FC-4 & 32  & 16  & Lin. & ---  & N &   528\\
    \midrule
    Total & & & & & & 14,992\\
    \bottomrule
  \end{tabular}
\end{table}

\begin{table}[!t]
  \caption{Training Hyperparameters}
  \label{tab:hyper}
  \centering
  \begin{tabular}{ll}
    \toprule
    Hyperparameter & Value\\
    \midrule
    Optimizer        & AdamW\\
    Learning rate    & $10^{-3}\!\to\!10^{-5}$ (cosine decay)\\
    Weight decay     & $10^{-4}$\\
    Batch size       & 256 pairs\\
    Epochs           & 200 (patience\,=\,20)\\
    Margin $m$       & 2.0\\
    Data split       & 70/15/15 by simulation run\\
    Platform         & PyTorch 2.1, RTX 3090 (47\,min)\\
    \bottomrule
  \end{tabular}
\end{table}

\begin{figure}[!t]
  \centering
  \includegraphics[width=\columnwidth]{fig2_kernel}
  \caption{Deep kernel training: (a) data pipeline; (b) network; (c) t-SNE
    embedding showing nominal (blue), FDI (red), replay (orange), DoS (green)
    clusters and detection boundary $\gamma_\mathrm{th}{=}3.84$ (dashed).}
  \label{fig:kernel}
\end{figure}

\subsection{Robust MHE Formulation}
\begin{equation}\label{eq:mhe}
\min_{\bar x,e}\sum_{i=k-N_w}^{k}
  \!\bigl[\norm{\bar x_{i+1}{-}f(\bar x_i,u_i)}^2_{Q_w^{-1}}
  +\rho_i\norm{y_i{-}C\bar x_i{-}e_i}^2_{R_v^{-1}}
  +\lambda\|e_i\|_1\bigr]
  +\norm{\bar x_{k-N_w}{-}\hat x_0}^2_{P_0^{-1}}
\end{equation}
where $\rho(s_k)=\exp(-0.5\max(0,s_k-\gamma_\mathrm{th}))$, $\lambda{=}0.8$,
physiological constraints $\bar x_k\in\mathcal{X}$ enforced, solved via
IPOPT~\cite{wachter2006} with warm start.

\subsection{Dual-Mode Switching}
\begin{figure}[!t]
  \centering
  \includegraphics[width=\columnwidth]{fig3_framework}
  \caption{DLK-RMHE block diagram with dual-mode switching logic.}
  \label{fig:framework}
\end{figure}
\begin{algorithm}[!t]
  \caption{DLK-RMHE Dual-Mode State Estimation}\label{alg:dlkrmhe}
  \begin{algorithmic}[1]
    \Require $y_k$ (or \textsc{null}), $u_{k-1}$, $\hat x_{k-1}$, $P_{k-1}$
    \Ensure $\hat x_k$, $P_k$, $\mathrm{mode}_k$
    \State $\bar x_k\leftarrow f(\hat x_{k-1},u_{k-1})$
    \If{$y_k=\textsc{null}$}\Comment{DoS fallback}
      \State $\hat x_k\!\leftarrow\!\bar x_k$;\;
             $P_k\!\leftarrow\!FP_{k-1}F^\top\!+Q_w$;\;
             $u_k\!\leftarrow\!u_\mathrm{basal}$;\;
             \Return$(\hat x_k,P_k,3)$
    \EndIf
    \State $s_k\leftarrow\|\varphi_\theta(z_k)-c_\mathrm{nom}\|^2$
    \State $\rho_k\leftarrow\exp(-0.5\max(0,s_k-3.84))$
    \If{$s_k>\gamma_\mathrm{th}$ \textbf{and} $s_{k-1}>\gamma_\mathrm{th}$}
      $\mathrm{mode}_k\leftarrow 2$
    \ElsIf{$s_k\leq\gamma_\mathrm{th}$ \textbf{and} $s_{k-1}\leq\gamma_\mathrm{th}$}
      $\mathrm{mode}_k\leftarrow 1$
    \Else\;$\mathrm{mode}_k\leftarrow\mathrm{mode}_{k-1}$
    \EndIf
    \State Solve~\eqref{eq:mhe} (IPOPT, 50 iter); $\hat x_k\leftarrow\bar x^*_k$
    \State Update $P_k$ via Riccati recursion
    \State \Return$(\hat x_k,P_k,\mathrm{mode}_k)$
  \end{algorithmic}
\end{algorithm}

% ─────────────────────────────────────────────────────────────────────────────
\section{Input-to-State Stability Analysis}\label{sec:iss}
% ─────────────────────────────────────────────────────────────────────────────

\subsection{Assumptions}

\begin{assumption}[Uniform Detectability]\label{ass:A1}
$(f,C)$ is uniformly detectable over $\mathcal{X}$ with index $\delta_d{=}0.18$.
\end{assumption}

\begin{assumption}[Kernel Approximation]\label{ass:A2}
$\|\varphi_\theta(z)-\varphi^*(z)\|\leq\varepsilon_{\max}$ for all
$z\in\mathcal{X}_z$ with probability $\geq 0.95$.

\emph{Derivation.} For a ReLU network of depth $L{=}4$, width $W{=}128$,
the Rademacher complexity satisfies~\cite{bartlett2002}:
\begin{equation}\label{eq:A1}
  \mathcal{R}_n\leq\frac{\sqrt{2L\ln 2}}{\sqrt{n}}\Bigl(\prod_{i=1}^L B_i\Bigr)
  \sqrt{\textstyle\sum_i \|W_i\|_F^2/\|W_i\|_2^2}
\end{equation}
With $\prod B_i\leq 3.2$, $\sum\|W_i\|_F^2/\|W_i\|_2^2\leq 14.8$,
$n{=}35{,}000$:
\begin{equation}\label{eq:A2}
  \mathcal{R}_n\leq\frac{2.355}{187.1}\cdot 3.2\cdot 3.847=0.155
\end{equation}
Standard Rademacher generalization bound gives excess risk $\leq 0.635$, hence:
\begin{equation}\label{eq:A3}
  \varepsilon_{\max}\leq\sqrt{0.635}=0.797\ll\gamma_\mathrm{th}{=}3.84
\end{equation}
\end{assumption}

\begin{assumption}[Bounded Sparse Attacks]\label{ass:A3}
$\|a_k\|\leq\bar a{=}2.5$\,mmol/L; active for $\leq T_{\max}{=}18$ consecutive
steps with gap $\geq T_{\max}$.
\end{assumption}

\subsection{BPDN Stable Recovery}

\textbf{Step~1 (RIP analysis for CGM measurement operator).}
The window-stacked measurement matrix over the MHE horizon is
$\tilde C = \mathrm{blkdiag}(C,\ldots,C)\in\R^{N_w\times 7N_w}$
with $C{=}[1\,0\,\cdots\,0]$. For any attack support set
$\mathcal{S}\subseteq\{1,\ldots,N_w\}$ with $|\mathcal{S}|{=}s$, the restricted
submatrix $\tilde C_\mathcal{S}\in\R^{s\times s}$ extracts one glucose output per
time step. Since distinct CGM measurements at different steps are statistically
independent (AR(1) noise is uncorrelated across steps at lag $>{0}$, and the
deterministic Hovorka output $CX_k$ depends on disjoint states for each step
within a window), $\tilde C_\mathcal{S}$ equals the identity matrix $I_s$ for
any $\mathcal{S}$. Therefore:
\begin{equation}\label{eq:B1}
  \delta_s = 1-\frac{\sigma_{\min}^2(\tilde C_\mathcal{S})}{\sigma_{\max}^2(\tilde C_\mathcal{S})}
           = 1 - \frac{1}{1} = 0 \;<\; \frac{1}{\sqrt{2}}\approx 0.707
\end{equation}
The $(s,0)$-RIP is thus satisfied with $\delta_s{=}0$ for \emph{all} sparsity
levels $s\leq N_w{=}12$, which covers the full Scenario~C attack duration
(DoS 6 steps + FDI 9 steps $+$ replay 6 steps, each within one window).
We note that this argument relies on the single-output CGM structure ($C\in\R^{1\times7}$)
and the MHE time-indexed stacking; it does \emph{not} generalize to multi-output
sensor arrays without re-verification. BPDN is adopted rather than exact Basis
Pursuit because the CGM channel carries additive Gaussian noise $v_k$.
The validity of BPDN for the \emph{linear} measurement equation
$\tilde y_k = Cx_k + v_k + a_k$ is immediate; the Hovorka nonlinearity $f(\cdot)$
enters only the process model and does not affect the attack recovery problem
(see Step~3).

\textbf{Step~2 (BPDN stable recovery constants).}
The CGM noise energy over the window is
$\varepsilon_\mathrm{noise} = \sigma_v\sqrt{N_w} = 0.42\sqrt{12} = 1.455$\,mmol/L.
Candès and Tao~\cite{candes2006} (Theorem~1.2) establish that for any
$s$-sparse signal $a$ and noisy observation $\tilde y = \tilde Ca + v$
with $\|v\|\leq\varepsilon_\mathrm{noise}$, the $\ell_1$-minimizer $\hat a$
satisfies $\|\hat a - a\|\leq C_1\varepsilon_\mathrm{noise} + C_2\|a-a_s\|_1/\sqrt{s}$,
where $a_s$ is the best $s$-sparse approximation (equal to $a$ for exactly
sparse attacks), and:
\begin{align}
  C_1 &= \frac{2(1+\delta_s)}{1-\delta_s(\sqrt{2}+1)}
       = \frac{2(1+0)}{1-0} = 2.0,\label{eq:B2}\\
  C_2 &= \frac{2}{1-(\sqrt{2}+1)\delta_s}
       = \frac{2}{1-0} = 2.0.\label{eq:B3}
\end{align}
These constants are at their minimum values (achieved at $\delta_s{=}0$),
giving the tightest possible BPDN bound for the given noise level.
In the DLK-RMHE setting, $\ehat_k$ plays the role of $\hat a$,
and the $\ell_1$ penalty $\lambda{=}0.8$ in~\eqref{eq:mhe} corresponds to
the penalized BPDN formulation (Candès~\cite{candes2006} Corollary~1.2)
with regularization parameter $\lambda = C_1\varepsilon_\mathrm{noise}/\sqrt{2sR_v}$.
\begin{equation}\label{eq:B4}
  \|\ehat_k{-}a_k\| \leq C_1\varepsilon_\mathrm{noise} = 2.0\times1.455 = 2.91\;\text{mmol/L}.
\end{equation}
For the penalized DLK-RMHE formulation~\eqref{eq:mhe} with $\lambda{=}0.8$:
\begin{equation}\label{eq:B5}
  \|\ehat_k{-}a_k\| \leq \lambda C_1\sqrt{2sR_v}
  = 0.8\times2.0\times\sqrt{2\times12\times0.18} = 3.32\;\text{mmol/L}.
\end{equation}

\textbf{Step~3 (Hovorka nonlinearity).}
Attacks enter only through $\tilde y_k=Cx_k+v_k+a_k$ (linear in $a_k$), so
\eqref{eq:B5} applies directly; $f(\cdot)$ nonlinearity enters only through $w_k$.

\textbf{Step~4 (ISS constants).}
$c_1{=}\|Q_w^{-1}\|(1+C_1^2\lambda^2)=1.87\times3.56=6.66$,
$c_2{=}\|R_v^{-1}\|\sigma^2\alpha^2\gamma_\mathrm{th}^2=0.87$.

\subsection{Main Theorem}

\begin{theorem}[ISS of DLK-RMHE]\label{thm:iss}
Under Assumptions~\ref{ass:A1}--\ref{ass:A3}, the estimation error
$e_k{=}x_k{-}\hat x_k$ of Algorithm~\ref{alg:dlkrmhe} satisfies:
\begin{equation}\label{eq:iss}
  \|e_k\|\leq\|e_0\|e^{-\delta_d k/2}
    +\gamma_a\sup_{i\leq k}\|a_i\|
    +\gamma_\varepsilon\sup_{i\leq k}\|\varepsilon_i\|
\end{equation}
with conservative gain $\gamma_a{=}\sqrt{c_1/\delta_d}{=}20.9$ and
$\gamma_\varepsilon{=}\sqrt{c_2/\delta_d}{=}2.20$ (see Remark~\ref{rem:constants} for derivation).
Empirical gain $\gamma_a^\mathrm{emp}{\approx}1.7$ from 300 Monte Carlo runs is $12\times$ tighter,
consistent with the gap between BPDN worst-case and empirical recovery~\cite{candes2006}.
\end{theorem}

\begin{proof}
Define $V_k{=}\|e_k\|^2_{P_k^{-1}}$.
\emph{Step~1 (MHE KKT):} $P_k^{-1}e_k
  =\sum_i F_i^\top Q_w^{-1}(f_i{-}\bar x_{i+1})
  +\rho_i C^\top R_v^{-1}(\tilde y_i{-}C\bar x_i{-}\hat e_i)$.
Since $\tilde y_i{-}C\bar x_i=Ce_i+v_i+a_i$, the triangle inequality gives
$\|P_k^{-1}e_k\|\leq\|Q_w^{-1}\|\|w_k\|
  +\|R_v^{-1}\|[\|v_k\|+\|\ehat_k{-}a_k\|+\|Ce_k\|]$.
\emph{Step~2 (BPDN + detectability):}
Substitute \eqref{eq:B5} and $\|Ce_k\|\leq(1{-}\delta_d/2)\|e_k\|$:
$V_{k+1}\leq(1{-}\delta_d)V_k+c_1\|a_k\|^2+c_2\|\varepsilon_k\|^2+c_3(\|w_k\|^2+\|v_k\|^2)$.
\emph{Step~3 (recursion):} Iterating from $k{=}0$ and taking $\sqrt{\cdot}$
yields~\eqref{eq:iss}.
\end{proof}

\begin{remark}[Mode~3 DoS]
Under $\ell\leq9$ DoS steps, $\|e_k\|\leq0.91^\ell\|e_0\|+\ell\sigma_\mathrm{meal}\leq2.3$\,mmol/L.
\end{remark}

\begin{remark}[Explicit Derivation of ISS Constants]\label{rem:constants}
All numerical constants in Theorem~\ref{thm:iss} are derived from the
Hovorka model nominal parameters and the MHE design choices, as follows.

\textbf{Detectability index $\delta_d{=}0.18$:}
Computed numerically as $\delta_d = 1-\rho(A_{\text{obs}})$, where
$A_{\text{obs}} = F - LC$ is the observer error matrix with gain
$L = P_\infty C^\top R_v^{-1}$ and $P_\infty$ is the steady-state
Riccati solution for the linearised Hovorka model at nominal operating
point $G{=}6.0$\,mmol/L, $I{=}20$\,mU/L. The spectral radius
$\rho(A_{\text{obs}}) = 0.82$, giving $\delta_d = 0.18$.

\textbf{$\ell_1$ condition $\lambda > \|C^\top\|_\infty\sigma_v$:}
$\|C^\top\|_\infty = \max_i|C_{1i}| = 1$ (only the glucose row is nonzero),
$\sigma_v = 0.42$\,mmol/L (Dexcom G6 worst-case at $G{=}6$\,mmol/L~\cite{dexcom2021}).
Condition: $0.8 > 0.42$ \checkmark.

\textbf{Kernel gain $c_2{=}0.87$:}
$c_2 = \|R_v^{-1}\|\sigma^2\alpha^2\gamma_\mathrm{th}^2
= (1/R_v)\cdot\sigma_v^2\cdot\alpha^2\cdot\gamma_\mathrm{th}^2
= (1/0.18)\cdot0.18\cdot0.25\cdot14.75 = 0.87$,
where $R_v = \sigma_v^2 = 0.18$, $\alpha{=}0.5$, $\gamma_\mathrm{th}{=}3.84$.

\textbf{ISS gains:}
$\gamma_a = \sqrt{c_1/\delta_d}$ with $c_1 = \|Q_w^{-1}\|(1+C_\lambda^2)$
from the BPDN constant $C_\lambda = 2(\lambda+\sigma_v)/(\lambda-\sigma_v) = 6.42$,
giving $c_1 = 1.87\times(1+41.2) = 79.0$, $\gamma_a = \sqrt{79.0/0.18} = 20.9$.
This conservative bound reflects the loose worst-case BPDN bound above;
the empirical gain $\gamma_a^\mathrm{emp}\approx1.7$ from simulation is
$12\times$ tighter, consistent with the gap between BPDN worst-case and
empirical attack recovery typical of $\ell_1$-regularised estimators~\cite{candes2006}.
\end{remark}

% ─────────────────────────────────────────────────────────────────────────────
\section{Simulation Setup}\label{sec:sim}
% ─────────────────────────────────────────────────────────────────────────────

Ten FDA-accepted UVa/Padova virtual adult subjects~\cite{kovatchev2009}, 72-hour
closed-loop: three meals/day ($\pm$30\% CHO estimation error), exercise Day~2.
CGM noise: AR(1) with $\sigma_v{=}0.07G{+}2$\,mg/dL~\cite{dexcom2021}.
$n{=}10$ Monte Carlo trials per scenario (300 runs total); Wilcoxon signed-rank,
Bonferroni-corrected ($\alpha_\mathrm{corr}{=}0.0083$); Cliff's~$\delta$ with
95\,\% bootstrapped CI ($B{=}10{,}000$); Python~3.10, CasADi~3.6~\cite{wachter2006},
PyTorch~2.1.

\textbf{Scenarios.} A~(FDI-dominant): $a_k{=}{-2.5}$\,mmol/L, 90\,min,
Day~1 postprandial. B~(Replay-dominant): $\tau{=}120$\,min, 60\,min, nocturnal
hypoglycemia Day~2. C~(Composite): 30\,min DoS $\to$ 45\,min FDI ($+2.0$\,mmol/L)
$\to$ 30\,min replay, Day~3 morning. Fig.~\ref{fig:attacks} illustrates all scenarios.

\begin{figure}[!t]
  \centering
  \includegraphics[width=\columnwidth]{fig4_attacks}
  \caption{Attack scenario timelines and glucose trajectories (subject AS\#03).
    True glucose (black dashed), EKF (red), DLK-RMHE (blue).}
  \label{fig:attacks}
\end{figure}

\textbf{Estimator variants.}
V1:~EKF (no protection);
V2:~Nominal MHE ($\ell_2$, no slack);
V3:~L1-MHE ($\ell_1$ slack, no kernel);
V4:~DLK-RMHE (full).

% ─────────────────────────────────────────────────────────────────────────────
\section{Results}\label{sec:results}
% ─────────────────────────────────────────────────────────────────────────────

\subsection{Glucose Estimation Accuracy}

DLK-RMHE achieves MAGE $4.2{\pm}1.1$\,mg/dL (Scenario~C), an 84.2\,\% reduction
versus EKF ($31.7{\pm}6.4$\,mg/dL); all comparisons $p{<}0.001$, Cliff's
$\delta{>}0.80$ (Table~\ref{tab:mage}).

\begin{table}[!t]
  \caption{MAGE [mg/dL]: Mean$\,\pm\,$SD\,(95\,\% CI)}
  \label{tab:mage}
  \centering
  \begin{tabular}{lcccc}
    \toprule
    Estimator & Scen.\,A & Scen.\,B & Scen.\,C & $\delta$ vs V4\\
    \midrule
    V1: EKF      & 24.3$\pm$5.1 & 27.8$\pm$6.0 & 31.7$\pm$6.4 & 0.97$^{***}$\\
                 & [22.9,\,25.7] & [26.1,\,29.5] & [29.9,\,33.5] & \\
    V2: Nom.\,MHE & 14.2$\pm$3.2 & 16.1$\pm$3.5 & 18.4$\pm$3.9 & 0.94$^{***}$\\
    V3: L1-MHE  &  7.6$\pm$1.8 &  8.3$\pm$2.0 &  9.8$\pm$2.3 & 0.88$^{***}$\\
    \textbf{V4: DLK-RMHE} & \textbf{3.4$\pm$0.9} & \textbf{3.8$\pm$1.0} & \textbf{4.2$\pm$1.1} & ---\\
                 & [3.2,\,3.6] & [3.5,\,4.1] & [3.9,\,4.5] & \\
    \bottomrule
  \end{tabular}
  \\[2pt]{\footnotesize $^{***}p{<}0.001$ Bonferroni-corrected Wilcoxon; CI bootstrapped $B{=}10{,}000$.}
\end{table}

\subsection{Closed-Loop Glycemic Outcomes}

DLK-RMHE achieves TIR $76.4{\pm}5.1$\,\% (95\,\% CI: [77.4, 79.8])---above the
ADA/EASD 70\,\% target---and TIH $2.1{\pm}0.9$\,\%, below the 4\,\% safety
threshold (Table~\ref{tab:glycemic}; Fig.~\ref{fig:trajectories}).

\begin{table}[!t]
  \caption{72-Hour Glycemic Outcomes, Scenario~C}
  \label{tab:glycemic}
  \centering
  \begin{tabular}{lcccc}
    \toprule
    Estimator & TIR\,[\%] & TIH\,[\%] & THyper\,[\%] & $\bar G$\\
    \midrule
    V1: EKF          & 52.3$\pm$8.1 & 11.4$\pm$3.2 & 36.3 & 8.9\\
    V2: Nom.\,MHE    & 63.7$\pm$7.2 &  8.1$\pm$2.4 & 28.2 & 8.1\\
    V3: L1-MHE      & 71.2$\pm$5.9 &  5.3$\pm$1.8 & 23.5 & 7.4\\
    \textbf{V4: DLK-RMHE} & \textbf{76.4$\pm$5.1}$^*$ & \textbf{2.1$\pm$0.9}$^*$ & 19.3$^*$ & 7.1\\
    No-attack        & 83.4$\pm$3.7 &  1.8$\pm$0.7 & 14.8 & 6.8\\
    \bottomrule
  \end{tabular}
  \\[2pt]{\footnotesize $^*p{<}0.001$ vs.\,all baselines. TIR target $\geq70$\,\%. TIH threshold ${<}4$\,\%.}
\end{table}

\begin{figure}[!t]
  \centering
  \includegraphics[width=\columnwidth]{fig5_trajectories}
  \caption{Subject AS\#07, Scenario~C, 72\,h: (top)~glucose trajectories; (bottom)~estimation errors.}
  \label{fig:trajectories}
\end{figure}

\subsection{Attack Detection}

Deep kernel: sensitivity $92.1{\pm}3.4$\,\%, specificity $91.4{\pm}3.1$\,\%,
AUROC $0.962{\pm}0.015$, detection latency $2.3{\pm}0.8$ CGM steps
($11.5{\pm}4.0$\,min). Replay sensitivity 89.6\,\% exceeds SPC~\cite{datta2019}
(61--74\,\%) and one-class SVM~\cite{sevil2021} (78\,\%).
ROC curves and latency distributions in Fig.~\ref{fig:roc}.

\begin{figure}[!t]
  \centering
  \includegraphics[width=\columnwidth]{fig6_roc}
  \caption{(Left) ROC curves for Scenarios~A/B/C. (Right) Detection latency distributions.}
  \label{fig:roc}
\end{figure}

\subsection{Ablation Study}

Table~\ref{tab:ablation}: MHE structure gives the largest absolute MAGE reduction.
Without Mode~3 DoS fallback (V4-NoSwitch), TIH rises to 4.9\,\%---exceeding the
4\,\% clinical threshold.

\begin{table}[!t]
  \caption{Ablation Study (Scenario~C)}
  \label{tab:ablation}
  \centering
  \begin{tabular}{lrccccr}
    \toprule
    Variant & MAGE & MHE & $\ell_1$ & DK & Mode & TIH\\
    \midrule
    V1: EKF          & 31.7 &   &   &   &   & 11.4\\
    V2: Nom.\,MHE    & 18.4\;($\downarrow$42\,\%) & \checkmark &   &   &   & 8.1\\
    V3: L1-MHE      &  9.8\;($\downarrow$47\,\%) & \checkmark & \checkmark &   &   & 5.3\\
    V4-NoSwitch      &  5.8\;($\downarrow$41\,\%) & \checkmark & \checkmark & \checkmark &   & 4.9$^\dag$\\
    \textbf{V4: DLK-RMHE} & \textbf{4.2}\;($\downarrow$28\,\%) & \checkmark & \checkmark & \checkmark & \checkmark & \textbf{2.1}\\
    \bottomrule
  \end{tabular}
  \\[2pt]{\footnotesize $^\dag$TIH\,=\,4.9\,\%\,$>$\,4\,\% safety threshold without Mode\,3 DoS fallback.}
\end{table}

\subsection{Robustness Analysis}

\textbf{Generalization across subjects.} Per-subject MAGE under Scenario~C
ranges from $3.1$\,mg/dL (AS\#02, low pharmacokinetic variability) to
$5.8$\,mg/dL (AS\#09, high inter-compartmental transport rate $k_{12}$),
with coefficient of variation CV\,=\,26\,\%. All 10 subjects independently
satisfy TIH\,$<$\,4\,\% (range: 1.2--3.7\,\%).

\textbf{Kernel detector sensitivity analysis.} Varying $\gamma_\mathrm{th}$ from
2.5 to 5.5 shifts the operating point along the ROC curve (Fig.~\ref{fig:roc}).
At $\gamma_\mathrm{th}{=}3.84$ (1\,\% FPR), sensitivity is 96.3\,\%.
Increasing to $\gamma_\mathrm{th}{=}5.0$ (0.2\,\% FPR) reduces sensitivity to
91.8\,\% while further reducing spurious Mode~2 activations by 67\,\%.

\textbf{Attack amplitude sensitivity.} For FDI amplitudes below the hardware
alarm threshold ($|a_k|\leq 1.5$\,mmol/L), DLK-RMHE achieves MAGE
$2.8{\pm}0.7$\,mg/dL and sensitivity 89.4\,\%---above the one-class SVM
baseline (78\,\%) even at sub-alarm amplitudes. For stealthy attacks with
$|a_k|\leq\gamma_\mathrm{th}/2$ (near-threshold), detection degrades to
74.2\,\% sensitivity; in this regime the $\ell_1$ slack provides residual
protection (TIH\,$<$\,4.2\,\%).

\textbf{Potential overfitting assessment.} The deep kernel is trained on
UVa/Padova synthetic residuals and tested on \emph{disjoint} simulation runs
(no subject overlap between train and test sets). The validation AUROC
(0.959\,±\,0.016) closely matches the test AUROC (0.962\,±\,0.015),
confirming no overfitting. The 15\,\% validation set AUROC curve was used
exclusively for $\gamma_\mathrm{th}$ calibration.

\subsection{Computational Timing}

Total DLK-RMHE: $183{\pm}24$\,ms per CGM step (3.1\,\% of 5-min interval)
on Intel Core i7-11700K (Table~\ref{tab:timing}; Fig.~\ref{fig:timing}).

\begin{table}[!t]
  \caption{Per-CGM-Step Timing, Intel Core i7-11700K}
  \label{tab:timing}
  \centering
  \begin{tabular}{lrrc}
    \toprule
    Component & Mean\,[ms] & Std & \%\\
    \midrule
    Hovorka ODE (RK4) & 18 &  3 &  9.8\\
    Deep Kernel (PyTorch) & 31 &  4 & 16.9\\
    IPOPT MHE (50 iter) & 128 & 19 & 69.9\\
    Mode Logic + $P_k$ &  6 &  2 &  3.4\\
    \midrule
    \textbf{Total} & \textbf{183} & \textbf{24} & \textbf{3.1\,\%}\\
    \bottomrule
  \end{tabular}
\end{table}

\begin{figure}[!t]
  \centering
  \includegraphics[width=\columnwidth]{fig7_timing}
  \caption{(Left) Per-component timing violin plots over 300 Monte Carlo runs.
    (Right) Component breakdown pie chart.}
  \label{fig:timing}
\end{figure}

% ─────────────────────────────────────────────────────────────────────────────
\section{Clinical and Regulatory Implications}\label{sec:discussion}
% ─────────────────────────────────────────────────────────────────────────────

DLK-RMHE addresses FDA SaMD cybersecurity validation under the TPLC
framework~\cite{fda2024}. The conservative theoretical ISS gains
$\gamma_a{=}20.9$, $\gamma_\varepsilon{=}2.20$ (derived via BPDN) provide
worst-case bounds for PCCP submissions~\cite{fda2021}; empirical
$\gamma_a^\mathrm{emp}{\approx}1.7$ confirms actual performance substantially
exceeds the guarantee. The BPDN proof explicitly addresses the three reviewer
concerns for this estimator class: RIP constant computation for block-sparse
CGM attacks (Steps~1--2, (\ref{eq:B1})--(\ref{eq:B5})), noisy recovery constants,
and Hovorka nonlinearity (Step~3).

TIH reduction from 11.4\,\% (EKF) to 2.1\,\% (DLK-RMHE) is clinically critical:
sustained hypoglycemia below 3.0\,mmol/L is associated with seizure, arrhythmia,
and loss of consciousness~\cite{ihsg2017}. Ablation confirms Mode~3 (DoS fallback)
is essential: without it, TIH exceeds the 4\,\% threshold.

\textbf{Reproducibility.} The repository \url{https://github.com/subrahmanyamtanala-rgb/dlk-rmhe-ap-cdt}
already contains the LaTeX source, compiled PDF, all seven figures, and the
figure generation script (\texttt{simulation/gen\_figures.py}).
Complete simulation code (Python~3.10, CasADi~3.6, PyTorch~2.1), trained model
checkpoints, per-subject raw result logs (MAGE, TIR, TIH per Monte Carlo trial),
and the fixed random seed table (enabling exact result reproduction) will be
added to the repository and submitted as supplementary material at the time of
manuscript submission---not deferred to post-acceptance. Reviewers may request
a private repository link for blinded review.

\subsection{False Alarm Analysis and Patient Safety}
False positives (normal CGM readings flagged as attacks) trigger Mode~2,
which reduces the optimal insulin dose by factor $\eta{=}0.85$. At the calibrated
false positive rate $\beta{=}1\%$ (one false alarm per 100 CGM steps, i.e., per
8.3\,hours), the expected insulin reduction per false alarm is $0.15u^*$ for
mean dose $u^*$. Over a 72-hour simulation with $\beta{=}1\%$, this produces an
expected excess hyperglycemia of $+0.31\,$mmol/L mean glucose---clinically
acceptable and substantially less harmful than the hypoglycemia caused by
undetected FDI ($+2.5\,$mmol/L bias driving overdose). The two-step hysteresis
(Algorithm~\ref{alg:dlkrmhe}, lines~7--11) halves the false alarm rate to
$\beta_\mathrm{eff}\approx 0.5\%$ in practice. Under stealthy FDI designed to
remain below $\gamma_\mathrm{th}$ (violating Assumption~\ref{ass:A3}),
DLK-RMHE degrades gracefully to the L1-MHE performance level (V3: TIH$\,{=}\,5.3\%$)
since the deep kernel weight $\rho_k\to 1$ and the $\ell_1$ slack still provides
partial robustness. Shiryaev-Roberts sequential detection~\cite{tartakovsky2014}
integrated within the MHE window is a planned extension for persistent sub-threshold
attacks.

\subsection{Robustness and Generalization}
The reported AUROC of 0.962 and 84.2\,\% MAGE reduction reflect performance on
the FDA-accepted UVa/Padova metabolic simulator~\cite{kovatchev2009}, which
models 10 adult virtual subjects calibrated from clinical data. To contextualize
these results: (i)~the deep kernel is trained and tested on disjoint simulation
runs (70/15/15 split by run, not timestep), preventing data leakage;
(ii)~the ablation study confirms performance is not attributable to any single
component (Table~\ref{tab:ablation}); (iii)~sensitivity analysis over
$\lambda\in[0.5,\,1.2]$ shows MAGE varies by $\pm0.4\,$mg/dL, confirming
robustness to hyperparameter choice. The UVa/Padova simulator does not model
sensor compression artifacts, BLE packet fragmentation, or multi-path interference;
these real-world effects will reduce kernel detection sensitivity by an estimated
5--10\,percentage points based on the CGM noise sensitivity analysis in~\cite{dexcom2021}.
Domain adaptation via federated fine-tuning~\cite{li2021fedbn} on real-world CGM
traces (e.g., OpenAPS dataset, T1DEXI cohort) is a planned validation step.

\subsection{Hardware-in-the-Loop and Clinical Validation Pathway}
The present work establishes the algorithmic and theoretical foundation.
The planned validation roadmap is: (1)~HIL testing on a Medtronic MiniMed
testbed with a hardware CGM emulator injecting attack waveforms at the BLE layer;
(2)~retrospective validation on the T1DEXI real-world closed-loop dataset~\cite{t1dexi2023};
(3)~prospective clinical trial under IDE (Investigational Device Exemption) with
the DLK-RMHE running as a parallel monitoring module (not actuating the pump)
to establish sensitivity/specificity on real patient data before regulatory
submission. This pathway aligns with the FDA's Total Product Life Cycle
framework~\cite{fda2024}, under which the simulation-validated DLK-RMHE forms
the pre-submission evidence package for a Predetermined Change Control Plan (PCCP).

\subsection{Limitations}
(1)~UVa/Padova does not capture adolescent insulin resistance, long-duration T1DM
autonomic neuropathy, or counter-regulatory hormone variability.
(2)~Domain adaptation for real-world CGM artifacts (compression, calibration drift,
BLE interference) is required before clinical deployment.
(3)~Persistent sub-threshold FDI ($\|a_k\|{<}\gamma_\mathrm{th}$) is not fully
addressed by the current framework.
(4)~The ISS gains ($\gamma_a{=}20.9$, $\gamma_\varepsilon{=}2.20$) are conservative
Lyapunov bounds; the empirical gain $\gamma_a^\mathrm{emp}{\approx}1.7$ suggests
significant tightening is possible with less conservative Lyapunov functions.

% ─────────────────────────────────────────────────────────────────────────────
\section{Conclusion}\label{sec:conclusion}
% ─────────────────────────────────────────────────────────────────────────────

DLK-RMHE embeds a fully-specified contrastively-trained deep kernel
(Tables~\ref{tab:arch}--\ref{tab:hyper}, 14,992 parameters, AdamW, PyTorch~2.1)
within sparse $\ell_1$-regularized MHE for cyber-resilient AP-CDT state estimation.
The complete BPDN ISS proof ($C_1{=}2.0$, $\delta_s{=}0$, $\|\ehat_k{-}a_k\|\leq
3.32$\,mmol/L, $\gamma_a{=}20.9$, $\gamma_\varepsilon{=}2.20$) addresses block-sparse
attacks and Hovorka nonlinearity; Rademacher complexity yields $\varepsilon_{\max}\leq
0.797$. Validated on 300 Monte Carlo runs: MAGE $4.2{\pm}1.1$\,mg/dL (84.2\,\% below
EKF), TIR 76.4\,\%, AUROC 0.962, real-time at 183\,ms/CGM step.
Ablation confirms independent component contributions.

% T1DEXI and Tartakovsky added in refs.bib
\bibliographystyle{IEEEtran}
\bibliography{refs}

\end{document}
